\chapter{Communication Complexity — Private Coin Finite Message}
%%% generated by the blueprint pipeline from TCSlib.CommunicationComplexity.NewmanTheorem.PrivateCoinFiniteMessage
\section{Overview}

This module defines private-coin finite-message protocols as a thin wrapper around
deterministic finite-message protocols, where Alice's input type is augmented with her
private randomness space $\Omega_X$ and Bob's with $\Omega_Y$.  It provides
constructors for Alice and Bob sending steps, an execution function \texttt{rrun} that
separates public inputs from private coins, approximate-correctness predicates
\texttt{CommunicationComplexity.PrivateCoin.FiniteMessage.Protocol.ApproxSatisfies} and \texttt{CommunicationComplexity.PrivateCoin.FiniteMessage.Protocol.ApproxComputes}, and conversions to and from the
binary private-coin protocol type.

%%%%%%%%%%%%%%%%%%%%%%%%%%%%%%%%%%%%%%%%%%%%%%%%%%%%%%%%%%%%%%%%%%%%%%%%%%%%%%%
\section{Declarations}

\begin{definition}[Private-coin finite-message protocol]
\lean{CommunicationComplexity.PrivateCoin.FiniteMessage.Protocol}
\label{CommunicationComplexity.PrivateCoin.FiniteMessage.Protocol}
\leanok
\uses{CommunicationComplexity.Deterministic.FiniteMessage.Protocol}
A private-coin finite-message protocol over input types $X$, $Y$ with private
randomness types $\Omega_X$, $\Omega_Y$ and output type $\alpha$ is defined as a
deterministic finite-message protocol whose effective input for Alice is
$\Omega_X \times X$ and for Bob is $\Omega_Y \times Y$.
\end{definition}

\begin{definition}[Output node]
\lean{CommunicationComplexity.PrivateCoin.FiniteMessage.Protocol.output}
\label{CommunicationComplexity.PrivateCoin.FiniteMessage.Protocol.output}
\leanok
\uses{CommunicationComplexity.PrivateCoin.FiniteMessage.Protocol}
The leaf node of a private-coin finite-message protocol that produces the constant
output $a : \alpha$, constructed by delegating to the corresponding output constructor
of the underlying deterministic finite-message protocol.
\end{definition}

\begin{definition}[Alice's sending step]
\lean{CommunicationComplexity.PrivateCoin.FiniteMessage.Protocol.alice}
\label{CommunicationComplexity.PrivateCoin.FiniteMessage.Protocol.alice}
\leanok
\uses{CommunicationComplexity.PrivateCoin.FiniteMessage.Protocol}
Given a function $f : X \to \Omega_X \to \beta$ (Alice's message function depending on
her public input and private randomness) and a continuation $P : \beta \to
\mathsf{Protocol}$, constructs a protocol node at which Alice sends a $\beta$-valued
message and execution branches on it.
\end{definition}

\begin{definition}[Bob's sending step]
\lean{CommunicationComplexity.PrivateCoin.FiniteMessage.Protocol.bob}
\label{CommunicationComplexity.PrivateCoin.FiniteMessage.Protocol.bob}
\leanok
\uses{CommunicationComplexity.PrivateCoin.FiniteMessage.Protocol}
Given a function $f : Y \to \Omega_Y \to \beta$ (Bob's message function depending on
his public input and private randomness) and a continuation $P : \beta \to
\mathsf{Protocol}$, constructs a protocol node at which Bob sends a $\beta$-valued
message and execution branches on it.
\end{definition}

\begin{definition}[Protocol execution with explicit coins]
\lean{CommunicationComplexity.PrivateCoin.FiniteMessage.Protocol.rrun}
\label{CommunicationComplexity.PrivateCoin.FiniteMessage.Protocol.rrun}
\leanok
\uses{CommunicationComplexity.Deterministic.FiniteMessage.Protocol.run, CommunicationComplexity.PrivateCoin.FiniteMessage.Protocol}
$\texttt{rrun}\,p\,x\,y\,\omega_x\,\omega_y$ executes the protocol $p$ on public inputs
$x : X$, $y : Y$ with Alice's private coin $\omega_x : \Omega_X$ and Bob's private coin
$\omega_y : \Omega_Y$, returning an element of $\alpha$.  It is defined by running the
underlying deterministic protocol on the paired inputs $(\omega_x, x)$ and
$(\omega_y, y)$.
\end{definition}

\begin{theorem}[Explicit-coin execution as paired-input execution]
\lean{CommunicationComplexity.PrivateCoin.FiniteMessage.Protocol.rrun_eq}
\label{CommunicationComplexity.PrivateCoin.FiniteMessage.Protocol.rrun_eq}
\leanok
\uses{CommunicationComplexity.Deterministic.FiniteMessage.Protocol.run, CommunicationComplexity.PrivateCoin.FiniteMessage.Protocol, CommunicationComplexity.PrivateCoin.FiniteMessage.Protocol.rrun}
\difficulty{1}
Let $p$ be a private-coin finite-message protocol over input types $X$, $Y$, with
private-randomness types $\Omega_X$, $\Omega_Y$ and output type $\alpha$. For all inputs
$x : X$ and $y : Y$ and all private coins $\omega_x : \Omega_X$ and $\omega_y :
\Omega_Y$, executing $p$ with those explicit coins agrees with executing the underlying
deterministic finite-message protocol on the paired inputs $(\omega_x, x)$ for Alice and
$(\omega_y, y)$ for Bob; that is, both yield the same element of $\alpha$.
\end{theorem}

\begin{definition}[Approximate satisfaction of a predicate]
\lean{CommunicationComplexity.PrivateCoin.FiniteMessage.Protocol.ApproxSatisfies}
\label{CommunicationComplexity.PrivateCoin.FiniteMessage.Protocol.ApproxSatisfies}
\leanok
\uses{CommunicationComplexity.PrivateCoin.FiniteMessage.Protocol, CommunicationComplexity.PrivateCoin.FiniteMessage.Protocol.rrun}
A protocol $p$ \emph{$\varepsilon$-satisfies} a predicate $Q : X \to Y \to \alpha \to
\mathrm{Prop}$ if for every input pair $(x, y)$,
\[
  \mathrm{vol}\bigl(\{(\omega_x, \omega_y) \mid \neg Q\,x\,y\,(p.\texttt{rrun}\,x\,y\,\omega_x\,\omega_y)\}\bigr) \;\le\; \varepsilon.
\]
\end{definition}

\begin{definition}[Approximate computation of a function]
\lean{CommunicationComplexity.PrivateCoin.FiniteMessage.Protocol.ApproxComputes}
\label{CommunicationComplexity.PrivateCoin.FiniteMessage.Protocol.ApproxComputes}
\leanok
\uses{CommunicationComplexity.PrivateCoin.FiniteMessage.Protocol, CommunicationComplexity.PrivateCoin.FiniteMessage.Protocol.rrun}
A protocol $p$ \emph{$\varepsilon$-computes} a function $f : X \to Y \to \alpha$ if
for every input pair $(x, y)$,
\[
  \mathrm{vol}\bigl(\{(\omega_x, \omega_y) \mid p.\texttt{rrun}\,x\,y\,\omega_x\,\omega_y \ne f\,x\,y\}\bigr) \;\le\; \varepsilon.
\]
\end{definition}

\begin{theorem}[Approximate computation as approximate satisfaction of equality]
\lean{CommunicationComplexity.PrivateCoin.FiniteMessage.Protocol.ApproxComputes_eq_ApproxSatisfies}
\label{CommunicationComplexity.PrivateCoin.FiniteMessage.Protocol.ApproxComputes_eq_ApproxSatisfies}
\leanok
\uses{CommunicationComplexity.PrivateCoin.FiniteMessage.Protocol, CommunicationComplexity.PrivateCoin.FiniteMessage.Protocol.ApproxComputes, CommunicationComplexity.PrivateCoin.FiniteMessage.Protocol.ApproxSatisfies}
\difficulty{1}
Let $p$ be a private-coin finite-message protocol over input types $X$, $Y$ with
private-randomness types $\Omega_X$, $\Omega_Y$ and output type $\alpha$, each of
$\Omega_X$ and $\Omega_Y$ carrying a measure, let $f : X \times Y \to \alpha$ be a
target function, and let $\varepsilon \in \bbr$. Then the proposition that $p$
$\varepsilon$-computes $f$ coincides with the proposition that $p$
$\varepsilon$-satisfies the predicate $Q(x,y,a)$ given by $a = f(x,y)$; that is, the two
propositions are equal.
\end{theorem}

\begin{definition}[Conversion to binary private-coin protocol]
\lean{CommunicationComplexity.PrivateCoin.FiniteMessage.Protocol.toProtocol}
\label{CommunicationComplexity.PrivateCoin.FiniteMessage.Protocol.toProtocol}
\leanok
\uses{CommunicationComplexity.Deterministic.FiniteMessage.Protocol.toProtocol, CommunicationComplexity.PrivateCoin.FiniteMessage.Protocol, CommunicationComplexity.PrivateCoin.Protocol}
Converts a private-coin finite-message protocol into a binary-tree private-coin
protocol by delegating to
\texttt{CommunicationComplexity.Deterministic.FiniteMessage.Protocol.toProtocol}.
\end{definition}

\begin{theorem}[Execution is preserved under conversion to a binary private-coin protocol]
\lean{CommunicationComplexity.PrivateCoin.FiniteMessage.Protocol.toProtocol_rrun}
\label{CommunicationComplexity.PrivateCoin.FiniteMessage.Protocol.toProtocol_rrun}
\leanok
\uses{CommunicationComplexity.Deterministic.FiniteMessage.Protocol.toProtocol_run, CommunicationComplexity.PrivateCoin.FiniteMessage.Protocol, CommunicationComplexity.PrivateCoin.FiniteMessage.Protocol.rrun, CommunicationComplexity.PrivateCoin.FiniteMessage.Protocol.toProtocol, CommunicationComplexity.PrivateCoin.Protocol.rrun}
\difficulty{1}
Let $p$ be a private-coin finite-message protocol with private-randomness spaces
$\Omega_X$ and $\Omega_Y$, input spaces $X$ and $Y$, and output type $\alpha$, in which
each message is an element of an arbitrary finite nonempty type. Fix inputs $x \in X$
and $y \in Y$ together with private coins $\omega_x \in \Omega_X$ and $\omega_y \in
\Omega_Y$. Then converting $p$ to a binary private-coin protocol leaves its randomized
execution unchanged: running the converted protocol on $x, y$ with coins $\omega_x,
\omega_y$ returns the same element of $\alpha$ as running $p$ itself on the same inputs
and coins.
\end{theorem}

\begin{theorem}[Complexity is preserved by the binary conversion]
\lean{CommunicationComplexity.PrivateCoin.FiniteMessage.Protocol.toProtocol_complexity}
\label{CommunicationComplexity.PrivateCoin.FiniteMessage.Protocol.toProtocol_complexity}
\leanok
\uses{CommunicationComplexity.Deterministic.FiniteMessage.Protocol.complexity, CommunicationComplexity.Deterministic.FiniteMessage.Protocol.toProtocol_complexity, CommunicationComplexity.Deterministic.Protocol.complexity, CommunicationComplexity.PrivateCoin.FiniteMessage.Protocol, CommunicationComplexity.PrivateCoin.FiniteMessage.Protocol.toProtocol}
\difficulty{1}
Let $p$ be a private-coin finite-message protocol with private randomness spaces
$\Omega_X$ and $\Omega_Y$, input types $X$ and $Y$, and output type $\alpha$, in which
each message is an element of an arbitrary finite nonempty type $\beta$ costing $\lceil
\log_2 \abs{\beta} \rceil$ bits. Converting $p$ into the associated binary private-coin
protocol, in which every message is encoded as a string of bits, leaves the worst-case
communication complexity unchanged: the complexity of the converted binary protocol
equals the complexity of $p$.
\end{theorem}

\begin{definition}[Embedding of binary private-coin protocol]
\lean{CommunicationComplexity.PrivateCoin.FiniteMessage.Protocol.ofProtocol}
\label{CommunicationComplexity.PrivateCoin.FiniteMessage.Protocol.ofProtocol}
\leanok
\uses{CommunicationComplexity.Deterministic.FiniteMessage.Protocol.ofProtocol, CommunicationComplexity.PrivateCoin.FiniteMessage.Protocol, CommunicationComplexity.PrivateCoin.Protocol}
Embeds a binary private-coin protocol into a private-coin finite-message protocol by
delegating to \texttt{Deterministic.FiniteMessage.Protocol.ofProtocol}.
\end{definition}

\begin{theorem}[Embedding preserves randomized execution of private-coin protocols]
\lean{CommunicationComplexity.PrivateCoin.FiniteMessage.Protocol.ofProtocol_rrun}
\label{CommunicationComplexity.PrivateCoin.FiniteMessage.Protocol.ofProtocol_rrun}
\leanok
\uses{CommunicationComplexity.Deterministic.FiniteMessage.Protocol.ofProtocol_run, CommunicationComplexity.PrivateCoin.FiniteMessage.Protocol.ofProtocol, CommunicationComplexity.PrivateCoin.FiniteMessage.Protocol.rrun, CommunicationComplexity.PrivateCoin.Protocol, CommunicationComplexity.PrivateCoin.Protocol.rrun}
\difficulty{1}
The embedding of a binary private-coin protocol into the finite-message framework does
not change what the protocol computes. Let $p$ be a binary private-coin communication
protocol with input spaces $X$ and $Y$, private-randomness spaces $\Omega_X$ and
$\Omega_Y$, and output type $\alpha$, and let $\widehat{p}$ be the finite-message
private-coin protocol obtained by embedding $p$ (interpreting each Boolean message as an
element of a two-element message alphabet). Then for every pair of inputs $x \in X$, $y
\in Y$ and every pair of private coins $\omega_x \in \Omega_X$, $\omega_y \in \Omega_Y$,
the randomized execution of $\widehat{p}$ agrees with that of $p$:
\[
\widehat{p}\text{'s output on }(x, y, \omega_x, \omega_y) \;=\; p\text{'s output on }(x,
y, \omega_x, \omega_y).
\]
\end{theorem}

\begin{theorem}[Embedding a binary private-coin protocol preserves complexity]
\lean{CommunicationComplexity.PrivateCoin.FiniteMessage.Protocol.ofProtocol_complexity}
\label{CommunicationComplexity.PrivateCoin.FiniteMessage.Protocol.ofProtocol_complexity}
\leanok
\uses{CommunicationComplexity.Deterministic.FiniteMessage.Protocol.complexity, CommunicationComplexity.Deterministic.FiniteMessage.Protocol.ofProtocol_complexity, CommunicationComplexity.Deterministic.Protocol.complexity, CommunicationComplexity.PrivateCoin.FiniteMessage.Protocol.ofProtocol, CommunicationComplexity.PrivateCoin.Protocol}
\difficulty{1}
Let $p$ be a binary private-coin protocol with randomness spaces $\Omega_X$ and
$\Omega_Y$, input spaces $X$ and $Y$, and output type $\alpha$. Then the private-coin
finite-message protocol obtained by embedding $p$ has the same communication complexity
as $p$; that is, viewing each Boolean message of $p$ as an element of the two-element
message alphabet leaves the worst-case number of bits exchanged unchanged.
\end{theorem}

\begin{theorem}[Faithful finite-message simulation of private-coin protocols]
\lean{CommunicationComplexity.PrivateCoin.FiniteMessage.Protocol.ofProtocol_equiv}
\label{CommunicationComplexity.PrivateCoin.FiniteMessage.Protocol.ofProtocol_equiv}
\leanok
\uses{CommunicationComplexity.Deterministic.FiniteMessage.Protocol.complexity, CommunicationComplexity.Deterministic.Protocol.complexity, CommunicationComplexity.PrivateCoin.FiniteMessage.Protocol, CommunicationComplexity.PrivateCoin.FiniteMessage.Protocol.ofProtocol, CommunicationComplexity.PrivateCoin.FiniteMessage.Protocol.ofProtocol_complexity, CommunicationComplexity.PrivateCoin.FiniteMessage.Protocol.ofProtocol_rrun, CommunicationComplexity.PrivateCoin.FiniteMessage.Protocol.rrun, CommunicationComplexity.PrivateCoin.Protocol, CommunicationComplexity.PrivateCoin.Protocol.rrun}
\difficulty{1}
Let $p$ be a binary private-coin protocol with randomness spaces $\Omega_X$ and
$\Omega_Y$, input spaces $X$ and $Y$, and output type $\alpha$, in which Alice's and
Bob's messages are single bits. Then there exists a private-coin finite-message protocol
$P$ over the same spaces such that, for every pair of inputs $x \in X$, $y \in Y$ and
every pair of private coins $\omega_x \in \Omega_X$, $\omega_y \in \Omega_Y$, the two
protocols return the same output when run on $x, y$ with coins $\omega_x, \omega_y$, and
moreover $P$ has exactly the same communication complexity as $p$.
\end{theorem}
